\problem{P32}{A Sum Stopped by Its Observations}{First-step decomposition}
Let $X_i$ be independent $\operatorname{Unif}(0,1)$ variables, and fix $0\le c<1$. Stop at the first $X_i>c$, writing $N$ for the stopping index and $S=\sum_{i=1}^NX_i$. Find $\E S$, $\Var S$, and $\Cov(S,N)$. Explain why the random-sum variance formula for an independent number of terms is inapplicable.
\solution Put $q=1-c$. Then $N$ is geometric with mean $1/q$. Since $0\le S\le N$, all required second moments are finite. With $(S',N')$ an independent copy of the continuation process,
\[
S=X+\ind\{X\le c\}S',\qquad
N=1+\ind\{X\le c\}N'.
\]
Let $\mu=\E S$ and $m_2=\E S^2$. Using $\E X=1/2$, $\E X^2=1/3$, and $\E[X\ind\{X\le c\}]=c^2/2$, we obtain
\[
\mu=\frac12+c\mu,\qquad
m_2=\frac13+c^2\mu+cm_2.
\]
Solving these equations gives
\[
\boxed{\E S=\frac1{2q},\qquad
\Var S=\frac{6c^2-4c+1}{12q^2}.}
\]
For $h=\E(SN)$, multiply the two first-step identities and take expectations:
\[
h=\frac12+\frac{c^2}{2q}+c\mu+ch.
\]
Thus $h=(1+c^2)/(2q^2)$ and
\[
\boxed{\Cov(S,N)=h-\mu/q=\frac{c^2}{2q^2}.}
\]
\pitfall Before stopping, observations are at most $c$; the terminal observation exceeds $c$. Conditioning on $N$ therefore changes their laws. The independent-count random-sum formula loses this dependence, whereas first-step decomposition retains it.
